% =====================================================================
\chapter{Atención como agrupación ponderada}
\label{tema:1}
% =====================================================================

Este tema define la atención en su forma más general: una consulta se compara con un conjunto de claves, los parecidos se convierten en pesos mediante una softmax y la salida es la media de los valores asociados ponderada por esos pesos. Se demuestran las dos propiedades estructurales básicas de la operación, la pertenencia de la salida a la envolvente convexa de los valores y su comportamiento en los límites de temperatura, y se presentan sus dos antecedentes: la regresión por núcleos de Nadaraya y Watson y la atención aditiva de Bahdanau, Cho y Bengio. Se siguen las §11.1, §11.2 y §11.4 de~\cite{d2l}, la §3.1 de~\cite{bahdanau}, la §7.1 de~\cite{jm3} y la clase~5 de~\cite{cs224n}.

\section{Del diccionario a la búsqueda suave}

Un diccionario almacena pares $(k_j,v_j)$ y, ante una consulta $q$, devuelve el valor $v_j$ cuya clave coincide con $q$, o un error si ninguna coincide. La búsqueda es \emph{dura}: el resultado depende de $q$ de forma discontinua y no admite gradientes. La atención sustituye la coincidencia exacta por un grado de parecido y la elección de un valor por una media ponderada~\cite[§11.1]{d2l}.

\begin{definicion}[Agrupación por atención]\label{def:agrupacion}
Sean $q\in\R^{1\times d_k}$ una \emph{consulta}, $k_1,\dots,k_n\in\R^{1\times d_k}$ unas \emph{claves} y $v_1,\dots,v_n\in\R^{1\times d_v}$ los \emph{valores} asociados. Sea $a\colon\R^{1\times d_k}\times\R^{1\times d_k}\to\R$ una \emph{función de puntuación}. La \emph{agrupación por atención} de los valores respecto de $q$ es
\[
\Att\bigl(q;(k_j,v_j)_{j=1}^{n}\bigr)=\sum_{j=1}^{n}\alpha_j\,v_j,\qquad
\alpha_j=\frac{\exp a(q,k_j)}{\sum_{l=1}^{n}\exp a(q,k_l)},
\]
es decir, $(\alpha_1,\dots,\alpha_n)=\softmax\bigl(a(q,k_1),\dots,a(q,k_n)\bigr)$. Los números $\alpha_j$ se llaman \emph{pesos de atención}.
\end{definicion}

La función de puntuación más usada es el producto escalar, $a(q,k)=q\cdot k=qk\T$, que es grande cuando los dos vectores apuntan en direcciones parecidas y tienen norma grande (módulo~2). Las claves deciden \emph{cuánto} contribuye cada par; los valores, \emph{qué} se contribuye. Separar ambos papeles permite que el criterio por el que un elemento es encontrado no coincida con la información que entrega.

\begin{figure}[h]
\centering
\begin{tikzpicture}[font=\small,>={Stealth[length=2mm]},
  caja/.style={draw,thick,rounded corners=2pt,minimum height=7mm,inner sep=4pt},
  op/.style={draw,thick,circle,inner sep=1.5pt}]
  \node[caja] (x) {$x_i$};
  \node[caja,right=1.2cm of x,yshift=9mm] (q) {$q_i=x_i\WQ$};
  \node[caja,right=1.2cm of x] (k) {$k_i=x_i\WK$};
  \node[caja,right=1.2cm of x,yshift=-9mm] (v) {$v_i=x_i\WV$};
  \draw[->] (x.east) -- (q.west);
  \draw[->] (x.east) -- (k.west);
  \draw[->] (x.east) -- (v.west);
  \node[caja,right=1.3cm of k,yshift=4.5mm] (s) {$s_{ij}=q_i\cdot k_j$};
  \node[caja,right=0.8cm of s] (sm) {$\alpha_{ij}=\softmax_j(s_{ij})$};
  \node[caja,below=0.9cm of sm] (o) {$o_i=\sum_j\alpha_{ij}v_j$};
  \draw[->] (q.east) -- (s.west);
  \draw[->] (k.east) -- (s.west);
  \draw[->] (s) -- (sm);
  \draw[->] (sm) -- (o);
  \draw[->] (v.east) -- (o.west);
\end{tikzpicture}
\caption{Autoatención para la posición $i$. Tres proyecciones aprendidas asignan a cada token los papeles de consulta, clave y valor; los parecidos de $q_i$ con todas las claves se normalizan con una softmax y ponderan los valores.}
\label{fig:qkv}
\end{figure}

\begin{definicion}[Autoatención]\label{def:autoatencion}
Sean $x_1,\dots,x_n\in\R^{1\times d}$ los vectores de los tokens de una secuencia y $\WQ,\WK\in\R^{d\times d_k}$, $\WV\in\R^{d\times d_v}$ matrices de parámetros. La \emph{autoatención} asigna a cada posición $i$ la salida $o_i=\Att\bigl(q_i;(k_j,v_j)_{j=1}^{n}\bigr)$, con
\[
q_i=x_i\WQ,\qquad k_j=x_j\WK,\qquad v_j=x_j\WV .
\]
Si las consultas proceden de una secuencia y las claves y valores de otra, se habla de \emph{atención cruzada}.
\end{definicion}

\begin{observacion}
La \cref{fig:qkv} resume la definición. En un modelo de lenguaje del tipo \emph{decoder-only}, como los que se estudian en este curso, toda la atención es autoatención; la atención cruzada aparece en los modelos codificador-decodificador de traducción (\cref{sec:bahdanau}).
\end{observacion}

\section{Propiedades de la agrupación}

\begin{proposicion}[Envolvente convexa]\label{prop:convexa}
Para cualesquiera $q$, claves, valores y función de puntuación, la salida de la \cref{def:agrupacion} pertenece a la envolvente convexa $\conv\{v_1,\dots,v_n\}$. En particular, para cada coordenada $c$,
\[
\min_j v_j[c]\;\le\;\Att(q)[c]\;\le\;\max_j v_j[c].
\]
\end{proposicion}

\begin{proof}
Los pesos son exponenciales divididas por su suma, luego $\alpha_j>0$ y $\sum_j\alpha_j=1$. Una combinación lineal de $v_1,\dots,v_n$ con coeficientes no negativos que suman~1 es, por definición, una combinación convexa. Para la cota por coordenadas, sean $m=\min_jv_j[c]$ y $m'=\max_jv_j[c]$; multiplicando $m\le v_j[c]\le m'$ por $\alpha_j>0$ y sumando en $j$ se obtiene $m\sum_j\alpha_j\le\sum_j\alpha_jv_j[c]\le m'\sum_j\alpha_j$, y se usa $\sum_j\alpha_j=1$.
\end{proof}

La atención, por tanto, solo \emph{mezcla} información existente: no puede producir una coordenada mayor que la mayor de los valores. Por esa razón un bloque Transformer combina la atención con una red de alimentación hacia delante y con conexiones residuales (módulo~14).

Para estudiar la dependencia respecto de la escala de las puntuaciones se introduce un factor $\beta>0$, el \emph{inverso de la temperatura} del módulo~1: $\alpha_j(\beta)=\exp(\beta s_j)/\sum_l\exp(\beta s_l)$, con $s_j=a(q,k_j)$.

\begin{teorema}[Límites de temperatura]\label{teo:limites}
Sean $s_1,\dots,s_n$ las puntuaciones y supóngase que el máximo $s^{\ast}=\max_js_j$ se alcanza en un único índice $j^{\ast}$. Entonces
\[
\lim_{\beta\to+\infty}\Att_\beta(q)=v_{j^{\ast}},\qquad
\lim_{\beta\to 0^{+}}\Att_\beta(q)=\frac{1}{n}\sum_{j=1}^{n}v_j .
\]
\end{teorema}

\begin{proof}
Dividiendo numerador y denominador por $\exp(\beta s^{\ast})$,
\[
\alpha_j(\beta)=\frac{\exp\bigl(\beta(s_j-s^{\ast})\bigr)}{\sum_l\exp\bigl(\beta(s_l-s^{\ast})\bigr)} .
\]
Si $j\ne j^{\ast}$, entonces $s_j-s^{\ast}<0$ y $\exp(\beta(s_j-s^{\ast}))\to 0$ cuando $\beta\to+\infty$; el término $j^{\ast}$ vale $e^{0}=1$. El denominador tiende a~1, luego $\alpha_{j^{\ast}}\to 1$ y $\alpha_j\to 0$ para $j\ne j^{\ast}$. Cuando $\beta\to 0^{+}$, cada $\exp(\beta s_j)\to 1$ por continuidad, el denominador tiende a $n$ y $\alpha_j\to 1/n$. En ambos casos el límite de la salida se obtiene sumando, porque la suma es finita.
\end{proof}

La atención interpola, pues, entre el diccionario exacto (límite $\beta\to\infty$, que devuelve el valor de la clave más parecida) y la media simple de los valores ($\beta\to 0$). El caso con empates se trata en el \cref{eje:1-limites}.

\section{Regresión por núcleos de Nadaraya y Watson}

La idea de promediar respuestas con pesos que dependen del parecido entre la consulta y unos puntos de referencia es anterior al aprendizaje profundo. Zhang \emph{et al.}~\cite[§11.2]{d2l} la presentan a partir de la regresión no paramétrica.

\begin{definicion}[Estimador de Nadaraya--Watson]\label{def:nw}
Sean $(k_j,v_j)_{j=1}^{n}$ observaciones de entrada y salida y $\kappa(q,k)\ge 0$ un \emph{núcleo}. El estimador de Nadaraya--Watson en el punto $q$ es
\[
\hat f(q)=\sum_{j=1}^{n}\frac{\kappa(q,k_j)}{\sum_{l}\kappa(q,k_l)}\,v_j .
\]
Con el núcleo gaussiano, $\kappa(q,k)=\exp\bigl(-\tfrac12\norm{q-k}^{2}\bigr)$.
\end{definicion}

\begin{proposicion}\label{prop:gauss}
Con el núcleo gaussiano, el estimador de Nadaraya--Watson es la agrupación por atención con puntuación $a(q,k)=-\tfrac12\norm{q-k}^{2}$. Si además todas las claves tienen la misma norma, coincide con la agrupación con puntuación $a(q,k)=q\cdot k$.
\end{proposicion}

\begin{proof}
La primera afirmación es la \cref{def:agrupacion} con $\exp a(q,k)=\kappa(q,k)$. Para la segunda, se desarrolla $-\tfrac12\norm{q-k_j}^{2}=q\cdot k_j-\tfrac12\norm{q}^{2}-\tfrac12\norm{k_j}^{2}$. Si $\norm{k_j}=r$ para todo $j$, los dos últimos términos no dependen de $j$. La softmax es invariante al sumar una misma constante a todas sus entradas (módulo~1), luego los pesos coinciden con los de $\softmax(q\cdot k_1,\dots,q\cdot k_n)$.
\end{proof}

\section{La atención aditiva de Bahdanau}
\label{sec:bahdanau}

En un traductor codificador-decodificador con redes recurrentes, el codificador produce un vector $h_j$ por cada palabra de la frase de origen. En el modelo clásico, el decodificador solo recibía el último de ellos, un resumen de tamaño fijo de toda la frase; Bahdanau, Cho y Bengio propusieron que en cada paso el decodificador construyese su propio resumen como media ponderada de todos los $h_j$~\cite[§3.1]{bahdanau}.

\begin{definicion}[Atención aditiva]\label{def:aditiva}
Sean $h_1,\dots,h_{T}$ los vectores del codificador y $s_{i-1}$ el estado del decodificador antes de emitir la palabra $i$. El \emph{vector de contexto} es
\[
c_i=\sum_{j=1}^{T}\alpha_{ij}h_j,\qquad \alpha_{ij}=\frac{\exp e_{ij}}{\sum_{l}\exp e_{il}},\qquad e_{ij}=v_a\T\tanh\bigl(W_as_{i-1}+U_ah_j\bigr),
\]
donde $W_a$, $U_a$ y $v_a$ son parámetros aprendidos (los vectores de esta definición son columnas, como en el artículo original).
\end{definicion}

Es una agrupación por atención en la que la consulta es $s_{i-1}$, claves y valores son los $h_j$ y la puntuación es una red de una capa oculta~\cite[apéndice~A.1.2]{bahdanau}. Vaswani \emph{et al.} señalan que la puntuación por producto escalar tiene una complejidad teórica similar, pero es mucho más rápida en la práctica porque se reduce a productos de matrices optimizados~\cite[§3.2.1]{vaswani}; es la que se estudia en el resto del módulo.

\begin{lecturas}
\item Zhang \emph{et al.}~\cite{d2l}, §11.1 «Queries, Keys, and Values», §11.2 «Attention Pooling by Similarity» y §11.4 «The Bahdanau Attention Mechanism».
\item Bahdanau, Cho y Bengio~\cite{bahdanau}, §3.1 «Decoder: General Description» y apéndice~A.1.2 «Alignment Model».
\item Jurafsky y Martin~\cite{jm3}, §7.1 «Attention».
\item CS224N~\cite{cs224n}, clase~5, «Transformers».
\end{lecturas}

\section*{Ejercicios del tema 1}
\addcontentsline{toc}{section}{Ejercicios del tema 1}

\begin{ejercicio}{1}{eje:1-kv}
Explique la diferencia entre el papel de la clave $k_j$ y el del valor $v_j$ en la \cref{def:agrupacion}. ¿Qué capacidad perdería el modelo si se impusiera $\WK=\WV$ (con $d_k=d_v$)?
\end{ejercicio}

\begin{ejercicio}{2}{eje:1-escalar}
Una consulta $q=(1,1)$ se compara, con puntuación $a(q,k)=q\cdot k/\sqrt{2}$, con las claves $k_1=(1,0)$, $k_2=(0,1)$ y $k_3=(1,1)$, cuyos valores escalares son $10$, $20$ y $30$. Calcule los pesos de atención y la salida.
\end{ejercicio}

\begin{ejercicio}{2}{eje:1-aditiva}
En la \cref{def:aditiva}, tómense $W_a=U_a=I_2$, $v_a=(1,1)\T$, el estado $s_{i-1}=(1,-1)\T$ y los vectores del codificador $h_1=(1,0)\T$, $h_2=(0,1)\T$, $h_3=(1,1)\T$. Calcule las puntuaciones $e_{ij}$, los pesos $\alpha_{ij}$ y el vector de contexto $c_i$.
\end{ejercicio}

\begin{ejercicio}{2}{eje:1-copiar}
Se usa la puntuación sin escalar $a(q,k)=q\cdot k$ con las claves ortonormales $k_1=(1,0,0)$, $k_2=(0,1,0)$, $k_3=(0,0,1)$ y los valores $v_1=(2,0)$, $v_2=(0,4)$, $v_3=(6,6)$. (a) ¿Qué condición deben cumplir la consulta y las claves para que casi todo el peso recaiga en un único índice? Compruébelo con $q=10\,k_3$. (b) Encuentre una consulta cuya salida sea aproximadamente $\tfrac12(v_1+v_2)$ y calcule la salida. (c) Repita el cálculo con $q=k_1+k_2$ y comente la diferencia.
\fuente{adaptado de CS224N~\cite{cs224n}, tarea~3, pregunta~1, apartados (a) y (b) (datos numéricos añadidos).}
\end{ejercicio}

\begin{ejercicio}{3}{eje:1-limites}
Generalice el \cref{teo:limites} al caso con empates: si el máximo $s^{\ast}$ se alcanza en un conjunto $J^{\ast}$ de $r$ índices, demuestre que $\lim_{\beta\to\infty}\Att_\beta(q)=\tfrac1r\sum_{j\in J^{\ast}}v_j$. Deduzca de la \cref{prop:convexa} que, si todos los valores son iguales, la salida no depende de la consulta.
\end{ejercicio}
